% Compile with: uplatex file.tex && dvipdfmx file.dvi
%           or: latexmk -uplatex file.tex
\documentclass[dvipdfmx]{article}
\usepackage{amsmath, amssymb}
\usepackage{tikz}

\title{scenario\_tight\_binding\_bands}
\date{\today}

\begin{document}
\maketitle

Register the lattice hopping model $H=-t\sum_{\langle ij\rangle}c^\dagger_ic_j$ (hopping amplitude $t$ between neighboring sites) symbolically, without expanding the sum.

$H = -t \sum_{\langle i j \rangle} c'_i c_j$

\begin{equation}
  H = -t \sum_{\langle ij \rangle} c'_i c_j
  \label{eq:E1}
\end{equation}

Find the dispersion relation $\varepsilon(k)$ (the eigenvalue of the Bloch matrix; for a single sublattice, simply the Fourier-transformed sum) treating this model as a 1D chain.

Find the dispersion relation of H

\begin{equation}
  \varepsilon(k) = -2t \cos k
  \label{eq:E2}
\end{equation}

Find the dispersion relation for the same hopping model on a square lattice.

Find the dispersion relation of H on a square lattice

\begin{equation}
  \varepsilon(\mathbf{k}) = -2t \left( \cos k_x + \cos k_y \right)
  \label{eq:E3}
\end{equation}

On a triangular lattice, find the dispersion relation including next-nearest-neighbor hopping $-t'$ in addition to nearest-neighbor hopping (adding the contribution of the second shell of neighbors).

Find the dispersion relation of H on a triangular lattice including next-nearest-neighbor -t'

\begin{equation}
  \varepsilon(\mathbf{k}) = -2t \left( \cos k_x + \cos k_y + \cos(k_x + k_y) \right) -2t' \left( \cos(k_x - k_y) + \cos(2 k_x + k_y) + \cos(k_x + 2 k_y) \right)
  \label{eq:E4}
\end{equation}

Find the dispersion relation on a honeycomb lattice (the graphene structure). With two sublattices, this should give two bands $\varepsilon_\pm(k)$ that are gapless at the Dirac point.

Find the dispersion relation of H on a honeycomb lattice

\begin{equation}
  \varepsilon_\pm(\mathbf{k}) = \pm t \sqrt{ 3 + 2\cos k_1 + 2\cos k_2 + 2\cos(k_1 - k_2) }
  \label{eq:E5}
\end{equation}

Add a sublattice-staggered on-site energy (a mass term $\Delta$) to the honeycomb lattice and find the dispersion relation. This should open a gap of $2\Delta$ at the Dirac point (the h-BN type).

Find the dispersion relation of H on a honeycomb lattice including a mass $\Delta$

\begin{equation}
  \varepsilon_\pm(\mathbf{k}) = \pm \sqrt{ \Delta^2 + t^2 \left( 3 + 2\cos k_1 + 2\cos k_2 + 2\cos(k_1 - k_2) \right) }
  \label{eq:E6}
\end{equation}

Find the dispersion relation on a kagome lattice (3 sublattices). Although not bipartite, this should yield a flat band plus two dispersive bands.

Find the dispersion relation of H on a kagome lattice

\begin{equation}
  \varepsilon_{\mathrm{flat}} = 2t, \quad \varepsilon_\pm(\mathbf{k}) = -t \left( 1 \pm \sqrt{ 3 + 2\cos k_x + 2\cos k_y + 2\cos(k_x + k_y) } \right)
  \label{eq:E7}
\end{equation}

Find the dispersion relation on a Lieb lattice (a flat band is expected from the star structure).

Find the dispersion relation of H on a Lieb lattice

\begin{equation}
  \varepsilon_0 = 0, \quad \varepsilon_\pm(\mathbf{k}) = \pm t \sqrt{ 4 + 2\cos k_x + 2\cos k_y }
  \label{eq:E8}
\end{equation}

Without using a named lattice, find the dispersion relation of a crystal specified directly by an abstract list of bonds (sublattice pair, lattice-vector displacement, amplitude). This should be a special case with the same structure as the SSH model.

Find the dispersion relation of the crystal with the following bonds: (0,1,(0),v),(0,1,(-1),w)

\begin{equation}
  \tilde H(k) = \begin{pmatrix} 0 & v + w e^{-i k} \\ v + w e^{i k} & 0 \end{pmatrix}
  \label{eq:E9}
\end{equation}

\begin{equation}
  \varepsilon_\pm(k) = \pm \sqrt{ v^2 + w^2 + 2 v w \cos k }
  \label{eq:E10}
\end{equation}

\end{document}